\documentclass[oneside]{ou-report}
\citestyle{agu}

% -------------------- Packages --------------------


% -------------------- Title metadata --------------------
\title[Master’s Thesis]{Requirements Traceability for Scriptless GUI Testing}
\author{R.\ J.\ H.\ van de Wetering}

\begin{document}

% ==================== Title page ====================
\begin{titlepage}
    \begin{center}
        % Logo
        \begin{tikzpicture}[remember picture,overlay]
            \node[anchor=south, inner sep=0pt] at (current page.south) {\includegraphics[scale=0.7]{logo}};
        \end{tikzpicture}
        
        \vspace*{2\bigskipamount}

        % Title
        {\makeatletter
            \titlestyle\color{title}\Huge\@title
        \makeatother}
    
        % Optional subtitle
        {\makeatletter
            \ifx\@subtitle\undefined\else
                \bigskip
                \titlefont\titleshape\LARGE\@subtitle
            \fi
        \makeatother}
        
        \vspace{2\bigskipamount}
        by
        \vspace{2\bigskipamount}
        
        % Author
        {\makeatletter
            \titlefont\Large\bfseries\@author
        \makeatother}
        
        \vfill
        
        in partial fulfillment of the requirements for the degree of
        
        \vspace{2\bigskipamount}
        
        {\bfseries Master of Science}
        
        in Software Engineering
        
        \vspace{2\bigskipamount}
        
        at the Open University, faculty of Management, Science and Technology \\
        Master Software Engineering
        
        to be defended publicly on Day Month DD, YYYY at HH:00 PM.
        
        \vfill

        % Metadata
        \begin{tabular}{lll}
            Student number:     & 852337023 \\
            Course code:        & \textsc{IM}9906 \\
            Thesis committee:   & Prof.\ dr.\ T.\ E.\ J.\ Vos (chair),  & Open University \\
                                & Dr.\ F.\ P.\ Ric\'os (supervisor),    & Polytechnic University of Valencia \\
        \end{tabular}
    \end{center}
\end{titlepage}

% ==================== Front matter ====================
\frontmatter

\chapter{Acknowledgements}

I would like to thank Prof.\ dr.\ T.\ E.\ J.\ Vos and Dr.\ F.\ P.\ Ric\'os for their supervision, constructive feedback, and ongoing support throughout this research. Our biweekly meetings provided valuable guidance, critical discussion, and thoughtful feedback that strengthened both the direction and execution of this work.

I am also grateful to the TESTAR team for providing the research context and tooling that enabled this study.

\chapter{Abstract}

Always give a summary in English.

\tableofcontents

\listoffigures

\listoftables

% ==================== Main matter ====================
\mainmatter

% ====================================================
% Introduction
% ====================================================
\chapter{Introduction}
You can start your report by copying the introduction from your VAF. The advice is to rewrite the introduction \emph{after} you did write the remainder of the thesis.

In the introduction, you:

\begin{itemize}
	\item motivate why this research is important. You can bring forward motivation with respect to society, or show that your research is scientifically interesting; preferrably both
	\item give some background information
	\item describe the goal of your research
	\item give the reader an overview of the structure of the remainder of your thesis.
\end{itemize}

This sentence is only here to show you how to refer to a source~\citep{Dijkstra-1968}.

\begin{table}[h!tbp]
\begin{tabular}{l | r | r| c}
kolom 1 & kolom 2 & kolom 3 & kolom 4 \\
\hline
zon & maan & ster & meteoor\\
gras & graan & groen & grauw\\
\end{tabular}
\caption{Example table}
\label{table-example}
\end{table}

Table~\ref{table-example} shows how to include a table. Note that the first column is left-justified, the right column is centered, and the other two columns are right-justified (because of the \texttt{\{l | r | r | c\}}). More information: \url{https://en.wikibooks.org/wiki/LaTeX/Tables}. 

\texttt{[h!tb]} means: preferrably place the table \emph{h}ere, and if that is not possible, at the \emph{t}op of the page, at the \emph{b}ottom, or on a separate \emph{page}. The same positioning advice can be used in figures. Figure~\ref{fig-example} is an example.

\begin{figure}[h!tbp]
\includegraphics[scale=0.5]{LaTeX.png}
\caption{LaTeX}
\label{fig-example}
\end{figure}

The following chapters are an example of how you could structure your thesis. Do not hesitate to use a different structure!

% ====================================================
% Problem analysis
% ====================================================
\chapter{Problem analysis}
Sometimes, it is sufficient to explain the problem and the background of the problem in the introduction. In such cases, you do not need a separate chapter about the problem.

At other times, you can show enough about the problem in your related work chapter.

But often, it is worthwhile to analyse the problem in a separate chapter. You could, for instance, describe a simple example, and show how your example raises questions.

% ====================================================
% Research
% ====================================================
\chapter{Research}
Start by copying everything about your research questions from your VAF. Always rewrite this chapter after you have finished your research results. Often, you will find that you did answer slightly different questions than you thought out before.

\section{Research questions}
In the introduction, you described the overall goal of your research.
Here, you formulate your research questions, and you explain them.

Note that your research questions must be formulated in such a way that you will be able to give meaningful answers in the conclusions of your thesis. The results you have produced must contain the answers to these questions.

In most cases, one main research question with several subquestions works best.

\section{Research method}
Here, you describe the method that you used to find answers to your questions. 

In general, it works well when you describe the method you use for each subquestion.  Therefore, you could opt for an alternative structure, by having subsections for each question, with the method described there as well.

\section{Validation}
You should not only describe how you find answers to your research questions, but also how you validate your work: how you will (try to) prove that your answers are indeed ansers to your questions.

Again, you can explain that for each subquestion, or here, in a separate section.

% ====================================================
% Related Work
% ====================================================
\chapter{Related Work}
You should find out what other researchers have found out about the problem that you will work on. You do that by searching for scientific sources that relate to the problem that you will work on. Often, the chapter in which you wrote your findings is called `Related Work'. Of course, you can choose another title.

You may conclude this chapter with a subsection in which you show how what you intend to do differs from what has been researched.

Try to structure this chapter around \emph{subjects} (instead of naming the different authors or articles that you found sequentially). The idea is that you (and your readers) get a clear view on what is known, and what is still unknown with respect to the problem.

You can, again, start by copying the related work section from your VAF to this chapter. After having done your research (and often during your research) add what you find.

% ====================================================
% Results
% ====================================================
\chapter{Results}
You will have at least one chapter describing your results. Often, it is worthwhile to dedicate a chaoter tp each subquestion. At other times, the methods you used cover several subquestions. In such cases, you could add a chapter for each method.

% ====================================================
% Conclusions
% ====================================================
\chapter{Conclusions}

Here, you give the answers to your questions and show how these answers are based on your results, and how they are validated.

% ====================================================
% Discussion
% ====================================================
\chapter{Discussion}
Here, you discuss everything that might be questionable about your research. Try to be very honest, and discuss everything that makes that the results are not 100\% reliable.

You may also make the discussion a subsection in your Conclusions.

Check whether the best place for the discussion is before or after the Conclusions.

% ==================== Back matter ====================
\backmatter

% ====================================================
% Reflection
% ====================================================
\chapter{Reflection}

\bibliographystyle{plainnat}
\bibliography{report}

% ==================== Appendices ====================
\appendix

\end{document}
